\documentclass[11pt]{article}

\usepackage[utf8]{inputenc}
\usepackage{tabto}
\usepackage[fleqn]{amsmath}
\usepackage{amssymb}
\usepackage{eurosym}
\usepackage{graphicx}
\usepackage{geometry}
\usepackage{hyperref}
\hypersetup{
    colorlinks,
    citecolor=black,
    filecolor=black,
    linkcolor=black,
    urlcolor=black
}

\renewcommand{\contentsname}{Inhaltsverzeichnis}

\begin{document}

\begin{titlepage}
	\begin{center}
		\vspace*{1cm}

		\Huge
		\textbf{Tourney}\\
		Projektplan

		\vspace{0.5cm}
		\LARGE
		Softwarepraktikum\\
		\Large
		Wintersemester 2014/2015

		\vspace{1.5cm}

		\large
		\textbf{Jonas Auer (2860992)\\
				 Fabian Biester (2859084)\\
				 Jan Tagscherer (2893134)}

		\vfill

		\includegraphics[width=0.4\textwidth]{Logo.png}

		\vspace{1.5cm}

		\Large
		Universität Stuttgart\\
		28.10.2014
	\end{center}
\end{titlepage}

\newpage

\tableofcontents
\newpage

\section{Einleitung}

\subsection{Zweck der Software}

Die Software soll die bisher händisch erfolgte Organisation von Turnieren übernehmen und vereinfachen. Dazu soll das Programm die Organisatoren durch den kompletten Prozess der Turnierplanung begleiten und sie bei den typischen Aufgaben wie der Teilnehmerverwaltung und der eigentlichen Durchführung der Turniere unterstützen.

\subsection{Motivation}

Momentan werden die Tuniere von der Turnierorganisation noch von Hand organisiert. Die große Menge an schriftlichen Notizen und Tabellen ist unübersichtlich und schwierig zu handhaben. Den zuständigen Organisatoren entsteht dadurch ein unnötiger, zeitraubender Aufwand, der gut durch eine entsprechende Software abgefangen werden kann und damit Zeit und Nerven spart. Diese Software würde zu einem übersichtlicheren und leichter zu handhabenden Prozess der Organisation von Turnieren beitragen.

\subsection{Projektüberblick}

\section{Formale Grundlagen}

\section{Leistungen der Vertragspartner}

\subsection{Lieferumfang}

Zum Lieferumfang des finalen Produkts gehören:

\begin{itemize}
	\item Projektplan
	\item Spezifikation
	\item Entwurf
	\item Systemtestprotokoll
	\item Das ausführbare Programm
	\item Der Quellcode des Programms
	\item Modultest
	\item Übderdeckungsprotokoll der Modultests
	\item Kurze Installations- und Startanleitung
	\item Handbuch
	\item Zeitabrechnung aller Teammitglieder
	\item Lizenz
\end{itemize}

\subsection{Leistungen des Auftraggebers}

\section{Anforderungen an die Umgebung}

\subsection{Infrastruktur}

An den Austragungsorten der Turniere ist eine dauerhafte Netzwerkanbindung nicht zu gewährleisten. Die Kommunikation zwischen den Geräten der beteiligten Organisatoren stellt sich dadurch als kompliziert heraus. Optimal wäre eine dauerhafte WLAN-Anbindung, allerdings wird aus Gründen der Zuverlässigkeit und Verfügbarkeit primär auf den Austausch von Daten über Massenspeichergeräten, wie zum Beispiel USB-Sticks, gesetzt.\\
Da viele verschiedene Geräte an dem Prozess beteiligt sind und diese zum Teil auf unterschiedliche Betriebssysteme setzen, müssen die Daten in einem platformunabhängigen Format bereitgestellt werden.

\subsection{Leistungen Dritter im Unternehmen}

\section{Entwicklungsprozess}

\subsection{Phasen der Entwicklung}

\subsection{Dokumentationsplan}

\subsection{Prüfung und Qualitätssicherung}

\section{Richtlinien für die Entwicklung}

\subsection{Konfigurationsmanagement}

\subsection{Design- und Programmierrichtlinien}

Der Fokus der Entwicklung liegt auf der intuitiven Bediunung des Programms. Die Software soll auch ohne das Handbuch gelesen zu haben eindeutig zu verstehen sein.\\
Im Hinblick auf den Programmcode wird Wert auf hohe Qualität gelegt. Dazu zählen die Einhaltung der [[Java Code Conventions hier eintragen]], sowie eine konsequente Dokumentation im Code mit Kommentaren und klarer Benennung von Bezeichnern in englischer Sprache.

\subsection{Eingesetzte Werkzeuge}

Zur Erstellung von Dokumenten und des eigentlichen Programms kommen eine Reihe von Werkzeugen zum Einsatz:

\begin{itemize}
	\item \textbf{Eclipse} - Entwicklungsumgebung
\end{itemize}

\section{Entwicklungsplan}

\subsection{Meilensteine}

\subsection{Arbeitspakete}

\subsection{Terminplan}

\subsubsection{Gantt-Diagramm}

\subsubsection{Termindrift-Diagramm}

\section{Risiken}

\subsection{Risiken und ihre Bewertung}

\subsection{Gegenmaßnahmen}

\section{Projektorganisation}

\subsection{Ansprechpartner und Beteiligte}

\end{document}
